% ===================================================================
% Section 3: A-, D-, E-optimality regularized portfolios
% ===================================================================
\section{A-, D-, and E-optimality regularized portfolios}
\label{sec:criteria}

\subsection{The unified objective}
Collecting the risk term, a ridge term, and the three design criteria
\eqref{eq:criteria}, the regularized portfolio solves
\begin{equation}
  \min_{\w\in\Wset}\ F(\w)
  \;:=\;
  \underbrace{\w^\top\Sighat\,\w}_{\text{risk}}
  \;+\;\underbrace{\lambda_2\lVert\w\rVert_2^2}_{\text{ridge}}
  \;+\;\underbrace{\gA\,\phiA(\M)}_{\text{A: average variance}}
  \;+\;\underbrace{\gD\,\phiD(\M)}_{\text{D: diversification}}
  \;+\;\underbrace{\gE\,\phiE(\M)}_{\text{E: robustness}} ,
  \label{eq:objective}
\end{equation}
with penalty strengths $\lambda_2,\gA,\gD,\gE\ge0$ and $\Wset$ as in
\eqref{eq:feasible}. Written out,
\begin{equation}
  F(\w)=\w^\top\Sighat\w+\lambda_2\lVert\w\rVert_2^2
  +\gA\tr\!\big(\M^{-1}\big)-\gD\log\det\M-\gE\lammin\!\big(\M\big).
  \label{eq:objective-explicit}
\end{equation}
In the square case $K=N$, Corollary~\ref{cor:KN} turns the middle two terms into
the familiar barriers, and \eqref{eq:objective-explicit} becomes
\begin{equation}
  F(\w)=\w^\top\Sighat\w+\lambda_2\lVert\w\rVert_2^2
  +\gA\sum_{i=1}^{N}\frac{G_{ii}}{w_i}
  -\gD\sum_{i=1}^{N}\log w_i
  -\gE\lammin\!\big(\M\big)
  \;+\;\text{const},
  \label{eq:objective-KN}
\end{equation}
which is the form used by our implementation and experiments. Setting all
penalties to zero recovers the GMV portfolio, and activating only the ridge
recovers the ridge-GMV portfolio; both statements are made precise in
Lemma~\ref{lem:reductions}.

\subsection{The three criteria and their economic meaning}

\paragraph{A-optimality: average precision of the factor premia.}
A-optimality minimizes the average variance of the parameter estimates,
$\phiA(M)=\tr(M^{-1})=\sum_{k}\lambda_k(M)^{-1}$. In the factor reading of
Section~\ref{sec:vi} this is the average estimation variance of the factor
premia $\theta=\mu_f$: the penalty pushes the portfolio toward the allocation
that pins down the \emph{average} factor premium most precisely. In the square
case it is the weighted inverse barrier \eqref{eq:Abarrier}, which diverges as
any $w_i\to0$ and, minimized alone over the simplex, is attained at
$w_i\propto\sqrt{G_{ii}}$ (Theorem~\ref{thm:limits}(ii)): a spread tilted by the
metric, rather than the uniform spread produced by D.

It deserves emphasis that A is a \emph{regularizer and not a risk minimizer}.
Minimizing average \emph{estimation} variance is not the same objective as
minimizing \emph{portfolio} variance, and the two can conflict; our experiments
document precisely this trade-off.

\paragraph{D-optimality: balanced factor coverage and diversification.}
D-optimality minimizes $\phiD(M)=-\log\det M=-\sum_k\log\lambda_k(M)$, and thus
maximizes the volume of the information ellipsoid. Because $\log\det$ is
maximized by spreading the spectrum, no factor direction is allowed to be
starved. In the square case the criterion separates into the log-barrier
$-\sum_i\log w_i$ \eqref{eq:Dbarrier}, which is maximized at $\w=\tfrac1N\one$
subject to the budget and pushes mass off the boundary: this is the spectral face
of diversification, and it is the link to the use of a D-optimality term directly
as a portfolio objective \cite{depaula2022}. In the rank-deficient case
$K<N$ the Cauchy--Binet form \eqref{eq:cb} shows that the criterion rewards
placing weight on subsets of assets whose loadings form a well-conditioned basis
of factor space (Remark~\ref{rem:scope}).

Note that the constant $-2\log|\det V|$ in \eqref{eq:Dbarrier} is independent of
$\w$; this is exactly why a criterion built on the raw covariance $\Sighat$
rather than on $\M$ would be inert. The criterion must act on the
weight-dependent information matrix.

\paragraph{E-optimality: worst-case robustness.}
E-optimality minimizes $\phiE(M)=-\lammin(M)$, that is, it maximizes the smallest
eigenvalue and so defends the least-informed direction in factor space. As a
variational object,
\begin{equation}
  \lammin(\M)=\min_{\lVert\uu\rVert_2=1}\uu^\top \M\,\uu
  =\min_{\lVert\uu\rVert_2=1}\ \sum_{i=1}^{N}w_i\,(\vv_i^\top\uu)^2 ,
  \label{eq:lammin-var}
\end{equation}
a pointwise minimum of functions \emph{linear} in $\w$, hence concave in $\w$.
Adding $-\gE\lammin(\M)$ is the finance image of the robust min--max portfolio
\cite{goldfarb2003}: it lifts the weakest eigenvalue, lowers the condition number
$\kappa(\M)=\lammax(\M)/\lammin(\M)$, and stabilizes the weights against
worst-case perturbations of the risk model. Where D is risk-neutral across
directions (it optimizes an average of log-eigenvalues), E is worst-case averse;
robust portfolio selection is worst-case averse, which is why E rather than D is
the natural design analogue of robustness.

\begin{remark}[Relation to ridge regularization]
Diagonal loading $\Sighat\mapsto\Sighat+\lambda I$ has an exact design
interpretation: with a Gaussian prior $\theta\sim\mathcal N(0,\tau^2I)$ the
Bayesian information matrix is $M(p)+(\sigma^2/\tau^2)I$, i.e.\ the information
matrix plus a constant ridge. Adding $\lambda I$ lifts every eigenvalue, so
$\lammin$ increases and $\kappa$ decreases. Ridge regularization is thus the
finance image of a Bayesian prior on the information matrix, and it acts on the
same spectral quantity that the E-criterion targets---but uniformly across
directions rather than adaptively on the weakest one. This is why we retain both
terms in \eqref{eq:objective}: they are complementary, a point the empirical
results in Section~\ref{sec:applications} bear out.
\end{remark}
